\beforeabstract
\prefacesection{Abstract}

This report covers Part A of COL334 Assignment 1, in which we measure and document the
network infrastructure of the IIT Delhi campus from the vantage point of our own devices.
Working as cartographers, we traced routes to a fixed set of targets and to a further set of
campus and external destinations, isolated the internal (\texttt{10.x.x.x}) hops that recur across
those traces, and used position, recurrence, and round-trip time to infer the likely role of each
hop in the campus topology. We surveyed the wireless layer from three locations, recording
every visible SSID, BSSID, band, channel, and signal strength, and used this to determine how
many physical access points are actually behind the broadcasts and how the campus's channel
plan is enforced. We then ran a set of ping-based delay experiments (fixed-size probes to five
targets at increasing distance, a packet-size sweep to isolate transmission delay, and, where
data was available, a queue-occupancy comparison between busy and quiet hours) to separate
propagation, transmission, and queuing delay in practice. Throughout, we have tried to keep the
line between what we directly measured and what we are inferring explicit, and to flag
alternative explanations we could not rule out rather than presenting a single tidy story.

\prefacesection{Acknowledgements}
We thank the COL334 teaching staff for the assignment design and the shared class map, and
IIT Delhi's network infrastructure for tolerating two days of low-rate probing.

\afterpreface \afterabstract
